\section{Evaluation}
\label{sec:evaluation}

The evaluation has three parts.  First, a fail-closed release evaluator binds
one proof, one public statement, and one production SBF binary to the theorem
and implementation artifacts.  Second, the exact binary is exercised on a
local Agave validator on both nullifier-account paths.  Third, a devnet
executor reproduces the account lifecycle and submits the same signed tag~60
wire that passed simulation.  The devnet result is admitted only after
finality and after a read-only evidence file has been committed.

\subsection{Release instance}
\label{sec:evaluation-release}

The release evaluator regenerated the default production SBF, replayed the
concrete proof in the host verifier, checked its three selector branches, and
reconciled 35 required gates.  All 35 gates passed and the failed-gate set was
empty.  The release-authorizing soundness value is the conservative
whole-ledger, three-branch floor of 100.16144938287457 bits.  The concrete
schedule is summarized in Table~\ref{tab:release-schedule}.

\begin{table}[t]
  \centering
  \small
  \caption{Concrete \Profile{} release schedule.}
  \label{tab:release-schedule}
  \begin{tabular}{l l}
    \toprule
    Parameter & Value \\
    \midrule
    profile / trace rows & 23 / \(2^{10}\) \\
    code rate & \(1/512\) \\
    query fibers & 18 \\
    batch work & 37 bits \\
    fold work & \(34,33,30,25\) bits \\
    final/query work & 32 bits \\
    selector candidates & 3 \\
    serialized selector & 0 \\
    accepted selector branches & \(0,1,2\) \\
    attempt cap & 17 \\
    release gates & 35 / 35 \\
    \bottomrule
  \end{tabular}
\end{table}

The content identities are as follows.
\begin{description}
  \item[Proof.] 66,367 bytes; SHA-256\\
  \texttt{f4e1e81f4a35b6b23f18430598ff98ec}\\
  \texttt{1f0db1146fabb4efd3c6715bcc847b53}.
  \item[Statement sidecar.] 667 bytes; SHA-256\\
  \texttt{976e9a7e001382025eaf81cfcb28ac609}\\
  \texttt{db966d4a9912511f54e2b702077b6de}.
  \item[Canonical public-input digest.] SHA-256\\
  \texttt{21d73e39be93112f986f52c7d683f2ab}\\
  \texttt{478890360a306af81110852ffb16a30a}.
  \item[Production SBF.] 915,656 bytes; SHA-256\\
  \texttt{da66a51b1f3ce95e907a87fca15fb9dc}\\
  \texttt{0cce66fd47646875ce2dff94879fd254}.
  \item[Release certificate.] SHA-256\\
  \texttt{b02eae619f685449a84095243f5d07be}\\
  \texttt{78f91ca7df67b435fb2c12aef8f39fcf}.
\end{description}

The statement has sequence zero, asset identifier 17, and fee 1.  The release
evaluator independently parsed all three post-final branches; all three met
the required ranks, so selector zero was both serialized and the least
\(\mathsf{Good}_{23}\) selector.  The proof and sidecar identities were required
to agree across the host-acceptance, local-mutation, and release artifacts.

\subsection{Local on-chain cost}
\label{sec:evaluation-local-cu}

Compute was measured by transaction simulation on
\texttt{solana-test-validator 2.3.0} (Agave source identifier
\texttt{a2e21dda}, feature set \texttt{3640012085}).  Each production
transaction requested the 1,400,000-CU limit and a 262,144-byte heap.  The
read-only tag~59 baseline and both tag~60 paths used the same uninstrumented
915,656-byte SBF\@.  Table~\ref{tab:production-cu} reports literal transaction
totals, rather than estimates from a traced diagnostic build.

\begin{table}[t]
  \centering
  \small
  \caption{Production-binary local compute consumption.}
  \label{tab:production-cu}
  \begin{tabular}{l r r r}
    \toprule
    Path & Tag~59 CU & Tag~60 CU & Headroom \\
    \midrule
    program-owned zero marker & 1,310,162 & 1,312,055 & 87,945 \\
    system-owned marker create & 1,310,162 & 1,314,386 & 85,614 \\
    \bottomrule
  \end{tabular}
\end{table}

The mutation wrapper added 1,893 CU when the marker was already
program-owned and zeroed, and 4,224 CU when the transaction created the
canonical PDA through the System Program.  The latter is the release maximum:
1,314,386 CU, or 93.8847\% of the requested limit.  Its remaining margin was
85,614 CU\@.

\subsection{Local functional and adversarial tests}
\label{sec:evaluation-local-tests}

The production-only test matrix compiled the explicit
\texttt{profile23-production} feature set and rejected its union with every
diagnostic and superseded Profile20--22 feature.  It then ran the exact
production tags on both supported marker paths.  The principal results are
listed in Table~\ref{tab:production-tests}.

\begin{table}[t]
  \centering
  \small
  \caption{Production-path state-transition tests.}
  \label{tab:production-tests}
  \begin{tabular}{p{0.56\linewidth} l}
    \toprule
    Test & Result \\
    \midrule
    mined proof accepted by tags 59 and 60 & pass \\
    committed unmined proof rejected by tags 59 and 60 & pass \\
    corrupted mined proof submitted as a transaction & rejected; exact rollback \\
    successful transition on each marker path & exact pool and marker images \\
    replay after a successful transition & rejected; no second mutation \\
    two-signer race on the system-owned marker path & exactly one commit \\
    upload and second finalization after tag~62 & rejected \\
    production tag~61 diagnostic path & unavailable \\
    \bottomrule
  \end{tabular}
\end{table}

The corruption cases were submitted, rather than only simulated, so their
post-transaction account images test rollback.  The successful cases checked
that the pool sequence advanced once, the output anchor replaced the current
anchor, and the marker contained the expected pool and nullifier.  The race
used two independently signed transactions targeting the same writable pool
and canonical nullifier PDA\@.

\subsection{Devnet method and evidence}
\label{sec:evaluation-devnet}

The read-only devnet preflight checked the release certificate, proof, sidecar,
and SBF identities against the canonical devnet genesis and deployed program:
\begin{center}
  \small
  genesis:\quad
  \texttt{EtWTRABZaYq6iMfeYKouRu166VU2xqa1wcaWoxPkrZBG}\\
  program:\quad
  \texttt{7Q2nGsPg8rbjdxKHK4jxTgEWLTyd9o1X4KMSjCieRmue}.\\
  ProgramData:\quad
  \texttt{cdRqe7MGCEJ2Z6iZfWuXtuymRiAhXtyfDgT47sKZr69}.
\end{center}
It observed that the deployed program already matched the 915,656-byte release
binary and that its ProgramData capacity was 925,760 bytes.  The preflight was
read only: it signed and submitted no transaction.

For the release proof, the fresh-account devnet schedule contains 109 setup
transactions followed by one tag~60 transaction.  The setup set consists of
pool creation, pool initialization, proof-account creation, proof-account
initialization, 104 chunk uploads, and proof finalization.  These transactions
are evidence for the prerequisite account lifecycle; they are not counted as
part of the one-transaction verification and state-transition result.

The preflight rent values were 1,224,960 lamports for the pool account,
463,083,600 lamports for the proof account, and 1,392,000 lamports for the
nullifier marker.  Their sum was 465,700,560 lamports.  The executor required
an additional 100,000,000-lamport fee reserve, giving a conservative
565,700,560-lamport readiness budget because deployment was not required.
This budget is a preflight upper allocation, not a measurement of transaction
fees.

Before final submission, the executor checks the complete uploaded account
image, finalizes it, and confirms that upload and second-finalization
simulations fail.  It snapshots the upgradeable program and ProgramData before
setup, immediately before simulation, and after finality, requiring byte and
capacity continuity.  The exact signed tag~60 transaction is simulated with
signature verification and without blockhash replacement.  The same serialized
bytes are then submitted; after an ambiguous RPC response, at most one
byte-identical retry is permitted.  The executor waits for the RPC
\texttt{finalized} state~\cite{SolanaRPC}, checks the pool and marker images,
and verifies that a duplicate simulation is rejected.

\subsubsection{Finalized result.}
\label{sec:evaluation-devnet-result}

The tag~60 transaction finalized successfully on devnet in slot
\ProfileDevnetSlot{} on
14 July 2026.  Its signature is
\begin{center}
  {\scriptsize\ttfamily
  3ofPbzRkqMEJZCM9vwKz96rLqRFtSg4d1GyqqVBEbogtwzmJodsWb2f7V4X83BLvuPXFsT6Yyf87PC1ZbLf1R7bx}.
\end{center}
The finalized record is available through the
\href{https://explorer.solana.com/tx/3ofPbzRkqMEJZCM9vwKz96rLqRFtSg4d1GyqqVBEbogtwzmJodsWb2f7V4X83BLvuPXFsT6Yyf87PC1ZbLf1R7bx?cluster=devnet}{Solana devnet explorer}.
The exact signed wire had SHA-256
\begin{center}
  {\scriptsize\ttfamily
  eae05cf339e825499277ef918451cac52595b1b4926c1ca3c081796e5e6be3d6}.
\end{center}
Its message had SHA-256
\begin{center}
  {\scriptsize\ttfamily
  80b1cdde261b624426a7ab012bd0c214151a3a469c91a6d6794c01cb097aa636}.
\end{center}

The executor submitted that same wire
after simulation.  Simulation and the finalized transaction both reported
\ProfileDevnetComputeUnits{} CU, and no identical-wire retry was used.  The
finalized RPC record
reported a 5,000-lamport transaction fee.

The finalized lifecycle contained 109 setup transactions followed by the one
tag~60 transaction.  The successful invocation resumed from the exact
\path{pool_initialized_proof_account_created_zeroed} checkpoint.  It
recovered the already-finalized pool creation, pool initialization, and
proof-account creation transactions by reconstructing their signed bytes from
the landed transactions and requiring byte-for-byte equality.  It also
required the exact initialized pool and zeroed proof-account images.  These
transactions were incorporated into the evidence and were not resent.

The 104 upload transactions each carried at most 640 proof bytes and consumed
643 CU.  They spanned finalized slots 476,227,548 through 476,231,535.  The
corresponding finalized block times were 15:15:53 and 15:40:37 UTC, a
24-minute-44-second upload interval.  This measurement reflects the executor's
serial policy of waiting for full finality after every chunk; it is not a
lower bound imposed by proof verification.

The proof account was finalized before verification.  Simulated upload and
second-finalization attempts were rejected without changing it.  Tag~60
advanced the pool sequence from zero to one, replaced its anchor, and created
the previously absent canonical nullifier marker.  A subsequent duplicate
simulation was rejected.  The principal account identities and data hashes
were:
\begin{description}
  \item[Proof account.]
  \texttt{HTSueSgUJxDbzjjtbeYJcLUzFajb2U3w1SNN25SWiknh}; finalized data
  SHA-256\\
  \texttt{b808f322a16110acb9c88e6f67ecefed}\\
  \texttt{e33c795e070a369e100bc371d4129f6d}.
  \item[Pool account.]
  \texttt{6QTNhChQjQpzL32gdtzR9RBdjZUSR3bQhNS3eTkgM6Fj}; data SHA-256 changed
  from\\
  \texttt{3aeb055bfa81d8aed6dc3a72bc4ae928}\\
  \texttt{3932841e3da893d4fc46de122bc67dd8}\\
  to\\
  \texttt{393633742ef61ea3d88f6a0f20832bae}\\
  \texttt{53781692888fcf2991fa8a977376c12a}.
  \item[Nullifier marker.]
  \texttt{AbWC9vKT7mxRNLeCnFrfi9HFsncF73XfG9TnbJS2jXHy}; data SHA-256\\
  \texttt{a276c72393f97b906563ab0c9416ed1c}\\
  \texttt{a943876eed56b77fb13180acef2870c2}.
\end{description}

The release-certificate identity remained
\texttt{b02eae619f685449a84095243f5d07be}\\
\texttt{78f91ca7df67b435fb2c12aef8f39fcf}.  The linked ProgramData identity was
\begin{center}
  \texttt{cdRqe7MGCEJ2Z6iZfWuXtuymRiAhXtyfDgT47sKZr69},\\
  deployment slot 476,224,357.
\end{center}
Program and ProgramData addresses,
complete account images, and data bytes were unchanged both immediately before
the final simulation and after finality.  No deployment transaction was
required.  The program remained upgradeable under the following authority:
\begin{center}
  \texttt{8AFmGcC2azHYGJz1VrDb8qJbCYdhVBTM5GE6TTSY3W2k}.
\end{center}

The resulting evidence file is 61,342 bytes, has mode 0444, and has SHA-256
\begin{center}
  \texttt{360e38fc5db3b644586c29e7a872203e}\\
  \texttt{8f9507c9ddef52add776fefb5d300275}.
\end{center}
These results are evidence for a devnet rehearsal; they are not mainnet
deployment evidence.

The completed evidence schema records, for every submitted transaction, its
label, signature, finalized slot, message hash, serialized-transaction hash,
reported CU, and identical-wire retry count.  Account evidence records the
address, lamports, owner, executable bit, data length, data hash, and
\[
  \operatorname{SHA256}(
    \mathsf{lamports}_{64}\mathbin\|\mathsf{owner}_{32}
    \mathbin\|\mathsf{executable}_{8}\mathbin\|\mathsf{data}).
\]
It also binds the release, proof, statement, selector, SBF, program snapshots,
proof-account seal, before-and-after pool state, nullifier state, and duplicate
rejection.  Execution first reserves the evidence pathname with an
\texttt{in\_progress\_no\_claim} marker.  Completion writes a new file,
flushes it, changes its mode to 0444, atomically renames it over the marker,
and flushes the parent directory.  An interrupted run therefore leaves an
explicit non-claim rather than a partial success record.

\subsection{Reproduction}
\label{sec:evaluation-reproduction}

The local production measurements and release certificate can be regenerated
from the repository root with the following commands.  The environment
variables must name the content-addressed proof and statement sidecar listed
in Section~\ref{sec:evaluation-release}.

\begin{verbatim}
export ASPIS_PROFILE23_PROOF=/absolute/path/to/profile23-proof.bin
export ASPIS_PROFILE23_STATEMENT=/absolute/path/to/profile23-statement.json

NO_DNA=1 cargo run --release -p aspis-xtask -- \
  stage2-atomic-profile23-acceptance
NO_DNA=1 cargo run --release -p aspis-xtask -- \
  stage2-atomic-profile23-mutation
NO_DNA=1 cargo run -q -p aspis-prover --example \
  profile23_soundness_epro_ledger
NO_DNA=1 cargo run --release -p aspis-xtask -- \
  stage2-profile23-one-transaction-release
\end{verbatim}

The final command removes and rebuilds the manifest-default production SBF
before comparing it with the explicit production-only binary and evaluating
all release gates.  Reproduction requires the proof and sidecar themselves;
their hashes in the release certificate are identifiers, not substitutes for
the files.  A publication artifact must therefore include those two files,
the generated SBF, the release certificate, and the finalized devnet evidence
under their recorded hashes.
